\documentclass[11pt,a4paper]{article}
\usepackage[T1]{fontenc}
\usepackage[utf8]{inputenc}
\IfFileExists{czech.ldf}{\usepackage[czech]{babel}}{\usepackage{babel}\babelprovide[main,import]{czech}}
\usepackage{lmodern,microtype,amsmath,amssymb,amsthm,booktabs,tabularx}
\usepackage[margin=24mm,headheight=15pt]{geometry}
\usepackage{fancyhdr}
\usepackage[colorlinks=true,linkcolor=black,urlcolor=blue]{hyperref}
\pagestyle{fancy}\fancyhf{}\fancyhead[L]{\small Kvantové vědomí: hypotéza, důkazy a otevřené otázky}
\fancyfoot[L]{\small Výzkumná syntéza pro Jana Mikulíka}\fancyfoot[R]{\thepage}
\setlength{\emergencystretch}{3em}\setlength{\parskip}{4pt}
\newtheorem{theorem}{Věta}
\newtheorem{proposition}[theorem]{Tvrzení}
\DeclareMathOperator{\Tr}{Tr}
\DeclareMathOperator{\TV}{TV}
\DeclareMathOperator{\diag}{diag}
\newcommand{\HH}{\mathrm H}
\newcommand{\II}{\mathrm I}
\begin{document}
\thispagestyle{empty}
{\LARGE\bfseries Kvantové vědomí, zaostřování a psychóza\par}
\medskip
{\large Výzkumná hypotéza, matematické důkazy a podklad pro odborný rozhovor\par}
\medskip
Zpracováno pro Jana Mikulíka na základě společné matematické práce.\\
4. listopadu 2026 \hfill Pracovní syntéza 1.0

\section*{Účel a status}
Tento dokument zachycuje společný matematický program a jeho možnou souvislost s popisem vnímání a psychózy. Rozlišuje tři druhy tvrzení:
\begin{enumerate}
\item \textbf{Matematické výsledky}, odvozené z vypsaných předpokladů.
\item \textbf{Empirické poznatky}, opřené o konkrétní publikované studie.
\item \textbf{Výzkumnou hypotézu}, která spojuje model s biologickým mechanismem a vyžaduje samostatné ověření.
\end{enumerate}
Cílem je připravit podklad pro rozhovor s lékařem a výzkumníkem. Text neobsahuje klinické vyšetření, diagnózu, rekonstrukci zdravotní historie ani doporučení změny léčby. Jan chce klinickou stránku probírat se svým lékařem; úlohou tohoto dokumentu je přesně vyložit výzkumný obsah.

\section*{Hlavní hypotéza}
\emph{V přesně vymezeném biologickém subsystému se kvantová dynamika může kauzálně podílet na procesech spojených s vědomým vnímáním, aktualizací přesvědčení a jednáním. Některé změny těchto procesů by se mohly projevit změnou citlivosti k podnětům, kontextu a zpětné vazbě.}

První část potřebuje fyzikální a biologické ověření. Druhá navíc vyžaduje definované měřitelné ukazatele a klinická data. Z matematických důkazů níže neplyne, že psychóza či schizofrenie jsou projevem kvantového vědomí.

\section*{Co je podloženo a co zůstává otevřené}
Dokazujeme korektnost instrumentového modelu, projektivní zaostření, klasickou simulovatelnost konečných záznamů, omezení entropie jako ukazatele správnosti a konkrétní komunikační separaci mezi jedním bitem a qubitem. Navrhujeme srovnání s klasickým modelem očekávání a evidence.

\textbf{Základní rozlišení:} subjektivní jistota, správnost vůči nezávislým údajům, kvantová koherence a klinický stav jsou různé veličiny.

Syntéza navazuje na pracovní texty \emph{Kvantová paměť jako možný mechanismus kauzální integrace vědomí}, \emph{Od obecné hypotézy k přesnému testu paměti} a \emph{Kvantová dynamika zaostřování a její klasická hranice}, všechny z 4.~listopadu 2026. Relevantní důkazy zde uvádíme samostatně. Netvrdíme ověřenou matematickou prioritu ani klinickou validaci.

\newpage
\section{Stav, měření, záznam a akce}
Nechť $\mathcal H=\mathbb C^d$ a $\rho_t=\rho_t^\dagger\succeq0$,
$\Tr\rho_t=1$. Historie je $h_t=(a_0,y_0,\ldots,a_{t-1},y_{t-1})$.
Politika volí akci $a_t=\pi(h_t)$. Relevantní paměť prostředí zahrnujeme do stavu; samotný redukovaný stav systému jinak nemusí pro předpověď stačit.

Každá akce vybírá konečný instrument:
\[
\mathcal M_y^{h,a}(\rho)=\sum_r K_{y,r}^{h,a}\rho(K_{y,r}^{h,a})^\dagger,
\qquad
\sum_{y,r}(K_{y,r}^{h,a})^\dagger K_{y,r}^{h,a}=I.
\]
Bornovo pravidlo a podmíněná aktualizace jsou
\[
p_y=\Tr\mathcal M_y^{h,a}(\rho_t),\qquad
\rho_{t+1}=\mathcal M_y^{h,a}(\rho_t)/p_y,\quad p_y>0.
\]
Bornovo pravidlo přijímáme jako fyzikální pravidlo standardního formalismu \cite{ibm}; neodvozujeme jeho fyzikální platnost z čisté algebry.

\begin{theorem}[Normalizace a pozitivita]
Čísla $p_y$ tvoří pravděpodobnostní zákon a definované posteriory jsou hustotní matice.
\end{theorem}
\begin{proof}
Pro libovolné $v$ je
$v^\dagger\mathcal M_y(\rho)v=\sum_r(K_{y,r}^\dagger v)^\dagger\rho(K_{y,r}^\dagger v)\geq0$.
Proto $p_y\geq0$. Cykličnost stopy a úplnost dávají
$\sum_y p_y=\Tr(\rho\sum_{y,r}K_{y,r}^\dagger K_{y,r})=1$.
Dělení kladnou stopou zachová pozitivitu a zajistí stopu jedna.
\end{proof}
Pro konečnou posloupnost výsledků definujeme
\[
P^\pi(\boldsymbol y)=
\Tr\!\left[\mathcal M_{y_{n-1}}^{h_{n-1},a_{n-1}}\circ\cdots
\circ\mathcal M_{y_0}^{h_0,a_0}(\rho_0)\right].
\]
Součet přes poslední výsledek zachovává stopu; iterace dává celkový součet jedna. Souřadnice $Y_t(\boldsymbol y)=y_t$ jsou skutečné náhodné veličiny. Na nulových větvích posterior nepotřebujeme.

\paragraph{Vůle a realizace.}
Vůli lze operačně modelovat politikou výběru akcí. Realizace v čase znamená konkrétní záznam interakcí. Tyto definice nepředstavují mechanismus uskutečňování libovolného přání. Změna pravděpodobností pod intervencí je odlišná od pouhé časové následnosti myšlenky a události.

\paragraph{Vně systému.}
Pro systém $S$ a prostředí $E$ lze psát
$\rho'_S=\Tr_E(U\rho_{SE}U^\dagger)$.
Prostředí je vně zvoleného subsystému, nikoli automaticky mimo fyzikální svět. Rovnice mu nepřisuzuje vědomí.

\newpage
\section{Zaostření a operátor logic--nonsense}
Pro projektivní rozlišení $\{P_i\}$, $\sum_iP_i=I$, definujeme
\[
q_i=\Tr(P_i\rho),\qquad U_{\mathcal P}(\rho)=\HH(q)
=-\sum_iq_i\log_2q_i,\quad 0\log_20:=0.
\]
Jde o neurčitost výsledku určitého měření, nikoli automaticky o neurčitost vědomí. Von Neumannova entropie $S(\rho)=-\Tr(\rho\log_2\rho)$ je jiná veličina.

\begin{theorem}[Projektivní zaostření]
Pro Lüdersův instrument a výsledek $i$ s $q_i>0$ platí
\[
\rho_i=P_i\rho P_i/q_i,\qquad \Tr(P_j\rho_i)=\delta_{ij},
\qquad U_{\mathcal P}(\rho_i)=0.
\]
Je-li $P_i=|e_i\rangle\langle e_i|$, pak $\rho_i=P_i$.
\end{theorem}
\begin{proof}
Použijeme $P_jP_i=\delta_{ij}P_i$ a cykličnost stopy.
Pro hodnost jedna
$P_i\rho P_i=\langle e_i|\rho|e_i\rangle P_i=q_iP_i$.
Dělení $q_i$ dává výsledek.
\end{proof}
Jde o opakované měření stejného rozlišení bez mezilehlé dynamiky. Při zapomenutí výsledku $\bar\rho=\sum_iP_i\rho P_i$ platí
$\Tr(P_j\bar\rho)=q_j$: původní entropie výsledku zůstane stejná.
Projektor vyšší hodnosti navíc nemusí odstranit vnitřní smíšenost stavu.

\begin{proposition}[Původní symetrizační operátor]
V ortonormální bázi $|L\rangle,|N\rangle$ nechť
$J|L\rangle=|N\rangle$, $J|N\rangle=|L\rangle$.
Pak $A=(I+J)/2$ je ortogonální projekce na
$\mathbb C|+\rangle$, kde $|+\rangle=(|L\rangle+|N\rangle)/\sqrt2$.
\end{proposition}
\begin{proof}
$J^\dagger=J$ a $J^2=I$ implikují $A^\dagger=A$ a $A^2=A$.
Dále
$A(a|L\rangle+b|N\rangle)=(a+b)|+\rangle/\sqrt2$
a $A|+\rangle=|+\rangle$, což určuje obraz.
\end{proof}
Pro $a+b=0$ je výstup nulový. Při fyzikální interpretaci je větev $A$ selektivní s pravděpodobností $\Tr(A\rho)$; druhá větev je $I-A$. Úspěch zvolené větve není zaručen.
Názvy logic a nonsense jsou označení vektorů, nikoli důkaz pravdivosti rozporu.

\paragraph{Co znamená nula.}
Konstantní veličina $Y\equiv1$ má $\HH(Y)=0$. Nulová entropie není nulová hodnota. Rovnost $\HH(Y\mid Y)=0$ znamená znalost zaznamenané veličiny, ne znalost celého světa.
Zobrazení $W(0)=W(1)=1$ sjednocuje výstupy, avšak v původní aritmetice stále $0\ne1$.

\newpage
\section{Shannonův most a klasické alternativy}
\begin{theorem}[Společný seed pro konečný strom]
Každý konečný řízený strom úplných historií s konečnými akcemi a výsledky a jádry $K(y\mid h,a)$ má jednu klasickou reprezentaci platnou pro všechny deterministické adaptivní politiky.
\end{theorem}
\begin{proof}
Pro každý pár $(h,a)$ vytvořme nezávislou souřadnici $S_{h,a}$ se zákonem $K(\cdot\mid h,a)$. Seed $S$ je celá konečná tabulka; dekodér vrací $S_{h,a}$. Dosažení $h$ závisí jen na souřadnicích předků, a proto neovlivní zákon $S_{h,a}$. Indukce po hloubce dává správnou pravděpodobnost každé větve. Stejné historie používají stejné souřadnice napříč politikami.
\end{proof}
Konstrukce nedokazuje minimálnost entropie seedu. Ta je samostatným tématem projektu \emph{finite-Shannon-theorem} \cite{repo}. Tabulka je konečná, ale může být velká. Tento dokument neaudituje všechny věty repozitáře.

\begin{theorem}[Diagonální instrumenty jsou klasicky realizovatelné]
Pokud ve společné pevné ortonormální bázi
\[
K_{y,r}^{h,a}=\sum_i k_{y,r,i}^{h,a}|e_i\rangle\langle e_i|,
\qquad
\ell_{h,a}(y\mid i)=\sum_r|k_{y,r,i}^{h,a}|^2,
\]
potom pro každou deterministickou adaptivní politiku
\[
P^\pi(\boldsymbol y)=\sum_iq_i\prod_t\ell_{h_t,a_t}(y_t\mid i),
\quad q_i=\langle e_i|\rho_0|e_i\rangle.
\]
\end{theorem}
\begin{proof}
Diagonální prvek výstupu $\mathcal M_y(\rho)$ je diagonální prvek $\rho$ násobený $\ell(y\mid i)$. Iterace dává součin a stopa součet.
Úplnost zaručuje $\sum_y\ell(y\mid i)=1$.
Pravá strana je klasický model: jednou vybereme skrytý index $i\sim q$ a další výsledky generujeme podle příslušných $\ell$.
\end{proof}
Výsledek platí i pro koherentní počáteční $\rho_0$: toto rozhraní koherenci nevidí. O celkové paměti řadiče nebo ceně náhodnosti zde netvrdíme malou horní mez.

K rozlišení mechanismů proto definujeme fyzikálně obhájenou třídu $\mathcal C(\mathcal R)$ s rozpočtem paměti, komunikace a času a zkoumáme
\[
d_{\mathcal R}(P)=\inf_{Q\in\mathcal C(\mathcal R)}
\sup_{\pi}\TV(P^\pi,Q^\pi),\qquad
\TV(P,Q)=\frac12\sum_z|P(z)-Q(z)|.
\]
Omezení se musí stanovit předem; nelze je vybrat až podle výsledků.
Sdílený seed sám neodstraňuje Bellovu lokalitu, pokud lokální výstupy závisejí jen na lokálním nastavení a seedu nezávislém na nastaveních.

\newpage
\section{Entropie, jistota a správnost}
V diagonálním modelu je posterior
$q_i^y=q_i\ell(y\mid i)/p_y$, kde $p_y=\sum_iq_i\ell(y\mid i)$.

\begin{theorem}[Průměrný informační zisk]
\[
\HH(q)-\sum_{y:p_y>0}p_y\HH(q^y)=\II(I;Y)\geq0.
\]
\end{theorem}
\begin{proof}
Pro společný zákon $p(i,y)=q_i\ell(y\mid i)$ je levá strana
$\HH(I)-\HH(I\mid Y)$.
To se rovná
\[
D_{\rm KL}(p_{IY}\Vert p_Ip_Y)
=\sum_{i,y:p(i,y)>0}p(i,y)\log_2\frac{p(i,y)}{q_ip_y}.
\]
Nezápornost plyne z $\ln u\leq u-1$ aplikované na poměr druhého zákona k prvnímu. Rovnost nastává při nezávislosti.
\end{proof}
Po jednotlivé větvi může entropie vzrůst:
\[
q=(0.9,0.1),\quad
K_y=\diag(\sqrt{0.1},\sqrt{0.9}),\quad
K_z=\diag(\sqrt{0.9},\sqrt{0.1}).
\]
Instrument je úplný. S pravděpodobností $p_y=0.18$ dostaneme
$q^y=(1/2,1/2)$; entropie vzroste z $0.4689956$ na $1$ bit.

\begin{proposition}[Jistota neurčuje správnost]
Správnost nejpravděpodobnější hypotézy není obecně funkcí samotné entropie subjektivního rozdělení.
\end{proposition}
\begin{proof}
Nechť nezávisle ověřený skutečný stav je $z_*=0$. Rozdělení
$\delta_0$ i $\delta_1$ mají entropii nula. Jejich nejpravděpodobnější hypotézy jsou však správná a chybná. Stejný vstup do funkce entropie nemůže určit dvě různé správnosti.
\end{proof}
Tento protipříklad není tvrzením o konkrétním člověku. Ukazuje, že zdravotní stav nelze bez dalších dat popsat jedním číslem nejistoty.

Také záleží na měření: $|0\rangle$ má $U_Z=0$, ale $U_X=1$.
Stavy
\[
|+\rangle\langle+|=\frac12\begin{pmatrix}1&1\\1&1\end{pmatrix},
\qquad I/2
\]
mají shodné $U_Z=1$, ale von Neumannovy entropie jsou nula a jeden bit.
Měření v bázi $|+\rangle,|-\rangle$ je odliší. To samo ještě nevylučuje obecnou klasickou simulaci.

V biologické interpretaci proto potřebujeme vedle nejistoty také kontrolovat přesnost, citlivost k novým údajům, kalibraci a kontext. Z matematické čistoty stavu nelze odvozovat zdravotní stav, hodnotu člověka nebo pravdivost jeho výkladu.

\newpage
\section{Očekávání, podnět a aktualizace přesvědčení}
Jako klasický srovnávací model zvolme konečné hypotézy $z$, kladný prior $q_t^-(z)$ a kladné věrohodnosti $L_t(y\mid z)$:
\[
q_t(z\mid y)=
\frac{[q_t^-(z)]^\alpha[L_t(y\mid z)]^\beta}
{\sum_u[q_t^-(u)]^\alpha[L_t(y\mid u)]^\beta},
\qquad \alpha,\beta>0.
\]
Pro $\alpha=\beta=1$ jde o Bayesovo pravidlo. Ostatní hodnoty jsou modelová volba, nikoli automaticky optimální inference nebo klinické parametry.

\begin{theorem}[Váha očekávání a evidence]
Pro dvě hypotézy platí
\[
\ln\frac{q_t(1\mid y)}{q_t(0\mid y)}
=\alpha\ln\frac{q_t^-(1)}{q_t^-(0)}
+\beta\ln\frac{L_t(y\mid1)}{L_t(y\mid0)}.
\]
\end{theorem}
\begin{proof}
V poměru posteriorů se zruší normalizace. Logaritmus součinu mocnin dává uvedený součet.
\end{proof}
Příklad: prior má poměr $9:1$, nový údaj poměr věrohodností $1:9$.
Pro $\alpha=\beta=1$ dostaneme $1:1$. Pro $\alpha=2,\beta=1$ zůstane $9:1$. To ilustruje převahu očekávání, ne diagnózu.

Je-li $\lambda=\ln[L_t(y\mid1)/L_t(y\mid0)]$,
pak při pevném prioru
\[
\frac{\partial}{\partial\lambda}\ln\frac{q_t(1\mid y)}{q_t(0\mid y)}=\beta.
\]
Na úlohách s nezávisle známou správnou odpovědí $z_t^*$ lze měřit například Brierovo skóre
\[
B=\frac1N\sum_{t=1}^N\sum_z
\left(q_t(z)-\mathbf 1_{\{z=z_t^*\}}\right)^2.
\]
Bez ověřeného cíle jej nelze poctivě vypočítat.
Počet odpovědí sám nerozliší změnu vnímání, rozhodovacího prahu nebo způsobu hlášení; tyto složky je třeba měřit odděleně.

\paragraph{Empirická opora, nikoli kvantový důkaz.}
Powers, Mathys a Corlett \cite{powers} zkoumali podmíněné sluchové vjemy pomocí učení, funkčního zobrazování a výpočetního modelu. Účastníci slyšící hlasy byli k efektu náchylnější. Studie podporuje zkoumání vlivu očekávání na vnímání. Neověřuje naše konkrétní exponenty, univerzální mechanismus psychózy ani kvantový původ vědomí.

Navržený kvantový model musí proti takové klasické alternativě přinést odlišnou, předem určenou a ověřitelnou předpověď. Pouhé přejmenování prioru na kvantový stav nestačí.

\newpage
\section{Podobnost s psychózou a její hranice}
Psychóza označuje soubor příznaků zahrnujících narušení vnímání, myšlení a rozlišování skutečnosti. Není synonymem schizofrenie a může se vyskytovat v různých souvislostech \cite{nimh}.
Halucinace a neobvyklá přesvědčení nejsou totožné jevy.

Prožitek může být skutečně prožívaný i bez odpovídajícího vnějšího podnětu. To samo neřeší jeho příčinu. Je nutné odlišit \emph{prožitek}, \emph{nezávisle zaznamenané okolnosti} a \emph{výklad prožitku}.
Sdělení člověka, že halucinace nemá, se má zachovat jako jeho sdělení; samo nerozhoduje o všech ostatních klinických otázkách.

\small
\noindent\begin{tabularx}{\textwidth}{@{}p{.22\textwidth}XX@{}}
\toprule
Zkoumaný jev & Kandidátní modelový popis & Co z podobnosti neplyne\\
\midrule
Vjem bez cílového podnětu & Očekávání, šum, práh odpovědi a atribuce zdroje & Vnější nebo kvantový původ\\
Silná jistota při nejasných údajích & Koncentrované $q$, váhy prioru a evidence & Správnost či diagnóza z jistoty\\
Obtížná oprava závěru & Citlivost k protidůkazu, dynamika učení & Vysvětlení celé psychózy jedním parametrem\\
Kontextová změna odpovědi & Stavová paměť či nekomutující instrumenty & Fyzikální kvantovost z pořadí otázek\\
Proměnlivý pocit autorství & Hypotézy o interním a externím zdroji & Kauzální ovládání děje z pocitu autorství\\
\bottomrule
\end{tabularx}
\normalsize

\paragraph{Hypotéza podobnosti.}
Některé subjektivní a behaviorální jevy lze modelovat změnami aktualizace, váhy očekávání, atribuce zdroje a korekce přesvědčení. Kvantově pravděpodobnostní model je jedna kandidátní formalizace. Jeho vhodnost musí určit porovnání s klasickými alternativami na stejných datech.

\paragraph{Silnější biologická hypotéza.}
Změny jsou způsobeny konkrétní změnou funkčně relevantní kvantové dynamiky. Pro toto tvrzení zde nemáme empirický důkaz. Vyžaduje nosič, měření a kauzální experiment.

Označení \emph{schizofrenie jako nezaostřeného kvantového vědomí} zatím nedává diagnostická kritéria ani měřitelný mechanismus. Protipříklady navíc ukazují, že nízká i vysoká entropie mají více významů. Testovatelná otázka zní: \emph{který aspekt aktualizace se mění, s jakým jevem souvisí a v čem kvantový model překoná klasický?}

Z analogie nelze odvodit zvláštní přístup k vnější realitě ani pravdivost konkrétního osobního výkladu. Text nepřiřazuje Janovi žádný řádek tabulky nebo hodnotu parametru.

\newpage
\section{Časové korelace a kontext}
\begin{proposition}[Časová korelace nevyžaduje kvantovou paměť]
Existuje klasický proces s jedním bitem vzájemné informace mezi libovolnými dvěma časovými záznamy.
\end{proposition}
\begin{proof}
Vyberme $B\sim\mathrm{Bernoulli}(1/2)$ a položme $Y_t=B$ pro každé $t$.
Pak $\HH(Y_t)=1$, $\HH(Y_t\mid Y_s)=0$, tedy $\II(Y_t;Y_s)=1$.
\end{proof}
Paměť, nemarkovovskost a kvantové provázání nejsou zaměnitelné pojmy. Pro fyzikální vícečasové procesy existují operační kritéria pracující s intervencemi \cite{pollock}.

\begin{proposition}[Pořadový efekt v kvantovém modelu]
Nechť $\rho=|0\rangle\langle0|$, $P=|0\rangle\langle0|$,
$Q=|+\rangle\langle+|$. Při dvou Lüdersových měřeních má dvojí odpověď ano pravděpodobnost $1/2$ v pořadí $P,Q$, ale $1/4$ v pořadí $Q,P$.
\end{proposition}
\begin{proof}
$\Tr(QP\rho PQ)=1\cdot|\langle+|0\rangle|^2=1/2$.
Naopak $\Tr(PQ\rho QP)=|\langle+|0\rangle|^2|\langle0|+\rangle|^2=1/4$.
\end{proof}
Pořadový efekt je tedy možný. Může jej ale reprodukovat i klasický model s aktualizací stavu a pamětí první otázky. Wang a spoluautoři \cite{wang} zkoumali konkrétní předpověď kvantové pravděpodobnosti pro pořadí otázek. Takové modelování samo neměří kvantový substrát mozku.

\paragraph{Akce a reakce.}
Klasická zpětná vazba už může mít tvar
\[
x_{t+1}\sim T(\cdot\mid x_t,a_t),\quad
y_t\sim O(\cdot\mid x_t),\quad a_t\sim\pi(\cdot\mid h_t).
\]
Vyjádřený názor může být fyzickou akcí a ovlivnit okolí.
Kauzální efekt je však rozdíl mezi rozděleními pod definovanými intervencemi, například
\[
P(Y\mid\operatorname{do}(A=a))-P(Y\mid\operatorname{do}(A=a')).
\]
Pouhá následnost myšlenky a události tento rozdíl neidentifikuje.

\paragraph{Kvantová kognice a fyzikální hypotéza.}
Hilbertovský model odpovědí a biologická kvantová paměť jsou rozdílné návrhy. Orch OR Hameroffa a Penrose \cite{orch} konkrétně předpokládá mikrotubuly a objektivní redukci. Náš obecný instrumentový model tento nosič ani redukci nepředpokládá. Citace návrhu není jeho experimentálním potvrzením.

\newpage
\section{Konkrétní test: jeden bit versus jeden qubit}
Příprava dostane rovnoměrný $x=(x_0,x_1)\in\{0,1\}^2$ a přenese jediný bit či qubit. Později čtenář dostane nezávisle vybraný $j\in\{0,1\}$ a hádá $x_j$. Nemá jiný přístup k původnímu vstupu. Sdílená náhoda je nezávislá na novém vstupu a otázce.

\begin{theorem}[Klasická mez]
Průměrný úspěch všech takových jednobitových klasických strategií je nejvýše $3/4$.
\end{theorem}
\begin{proof}
Při pevné náhodnosti je strategie deterministická. Dekodér má dva odpovědní vektory $v_0,v_1\in\{0,1\}^2$. Alespoň dva ze čtyř vstupů nejsou žádným z nich. Každý takový vstup způsobí alespoň jednu chybu přes dvě otázky. Z osmi párů je tedy nejvýše šest správných. Směsi mez nezvýší.
\end{proof}
Rovnoměrného úspěchu $3/4$ dosáhneme sdíleným spravedlivým indexem $k$: pošleme $x_k$, při $j=k$ odpovíme přesně, jinak náhodně.

Kvantová rodina
\[
\rho_x^{(\eta)}=\frac12\left[I+\frac{\eta}{\sqrt2}
\big((-1)^{x_0}\sigma_x+(-1)^{x_1}\sigma_z\big)\right],
\quad 0\leq\eta\leq1
\]
je pozitivní: délka Blochova vektoru je $\eta$. Zde $\sigma_x,\sigma_z$ jsou Pauliho matice. Měření $\sigma_x$ pro $j=0$ a $\sigma_z$ pro $j=1$ dává
\[
p_\eta=\Pr(b=x_j\mid x,j)=\frac12+\frac{\eta}{2\sqrt2}.
\]
Jde o známý quantum random access code \cite{rac}.

\begin{theorem}[Přesná vzdálenost]
Pro uvedenou klasickou třídu $\mathcal C_1$ platí
\[
\inf_{Q\in\mathcal C_1}\max_{x,j}\TV(P_\eta(\cdot\mid x,j),Q(\cdot\mid x,j))
=\max\{0,\eta/(2\sqrt2)-1/4\}.
\]
\end{theorem}
\begin{proof}
Maximální TV chyba $r$ znamená klasický úspěch alespoň $p_\eta-r$,
tedy $r\geq p_\eta-3/4$.
Je-li $p_\eta\geq3/4$, rovnoměrně optimální klasický model mez dosahuje.
Jinak jej smícháme s náhodnou odpovědí s vahou $4p_\eta-2$. V každém nastavení dostaneme úspěch $p_\eta$, což pro binární výstup určuje celý zákon.
\end{proof}
Dva klasické bity uloží celý vstup a rodinu simulují přesně. Překročení $3/4$ v lidské úloze proto neprokazuje kvantový mozek bez fyzikálně zajištěného jednobitového omezení a kontroly dalších kanálů. Benchmark sám netestuje vědomí, psychózu ani provázání s prostředím.

\newpage
\section{Výzkumný plán a podklad pro rozhovor}
\begin{enumerate}
\item \textbf{Operationalizace.} Určit měřený jev: odpověď, jistotu, chybu atribuce zdroje nebo změnu závěru. Každý potřebuje vlastní ukazatel.
\item \textbf{Alternativy.} Srovnat váženou klasickou inferenci, klasický stavový model a kvantově pravděpodobnostní model při srovnatelné složitosti.
\item \textbf{Identifikovatelnost.} Na simulovaných datech ověřit, zda lze parametry vůbec rekonstruovat a rozlišit.
\item \textbf{Predikce.} Předem stanovit testovací podmínky a hodnotit nové, odložené údaje. Shoda s trénovacími daty sama nestačí.
\item \textbf{Fyzikální nosič.} Pro biologickou kvantovou hypotézu určit subsystém, dobu koherence, přípravu, čtení a vazbu na neuronové procesy.
\item \textbf{Kauzální relevance.} Ověřit, že cílená změna navrženého mechanismu mění předem určenou funkci, a kontrolovat klasická vysvětlení.
\item \textbf{Klinická vazba.} Teprve pak zkoumat vztah ke konkrétním klinickým jevům, s odborným a etickým zajištěním a nezávislou replikací.
\end{enumerate}
Konkrétní varianta neobstojí, pokud její předpovědi nevycházejí na nových datech, potřebný nosič nevykazuje požadované vlastnosti nebo rozdíl vysvětluje povolená klasická alternativa. Neúspěch jedné varianty nevyvrací všechny myslitelné hypotézy.

\paragraph{Pro rozhovor s lékařem.}
Jan žádá, aby jeho dlouhodobá práce byla brána vážně a aby výzkumná hypotéza nebyla zaměněna s posouzením jeho osoby. Tento dokument nepřisuzuje Janovi diagnózu, parametry ani příčiny prožitků. Nenahrazuje jeho vlastní popis zkušenosti.

K rozhovoru lze připravit konkrétní příklady a otázky: Které prožitky jsou přítomné a které nejsou? Co je vjem, co interpretace a co matematická metafora? Jak se mění spánek, soustředění, zátěž, jistota a každodenní fungování? Co má nezávislé potvrzení? Které alternativní vysvětlení lze rozlišit pozorováním? Odpovědi nelze doplnit z rovnic.

\paragraph{Status důkazů a práce.}
Text obsahuje analytické důkazy abstraktních výsledků. Předchozí poznámky obsahují numerické kontroly instrumentů a enumeraci klasických strategií; ty nejsou klinická data. Tato syntéza není certifikát Leanu. Důkaz kvantového původu vědomí nebo psychózy zde nemáme.

Smyslem práce je převést intuici na podmíněná tvrzení, kontrolovatelné důkazy a rozlišující experimenty. Prožitek, jeho interpretace, matematická konzistence a biologická pravdivost mají každý vlastní způsob ověřování.

\newpage
\section*{Zdroje a jejich role}
Zdroje kontrolovány 4.~listopadu 2026. Jde o cílenou rešerši, nikoli systematický přehled celé literatury. Klinické souvislosti nejsou vydávány za důsledky matematických vět.

\begin{thebibliography}{9}
\bibitem{nimh}
National Institute of Mental Health.
\emph{Understanding Psychosis}. NIH Publication No.~23-MH-8110.
\url{https://www.nimh.nih.gov/health/publications/understanding-psychosis}.
Vymezení klinických pojmů.

\bibitem{powers}
A.~R. Powers, C. Mathys a P.~R. Corlett.
Pavlovian conditioning-induced hallucinations result from overweighting of perceptual priors.
\emph{Science} 357(6351), 596--600 (2017).
\url{https://doi.org/10.1126/science.aan3458}.
Empirická studie vnímání; nikoli test kvantového mechanismu.

\bibitem{wang}
Z. Wang, T. Solloway, R.~M. Shiffrin a J.~R. Busemeyer.
Context effects produced by question orders reveal quantum nature of human judgments.
\emph{PNAS} 111(26), 9431--9436 (2014).
\url{https://doi.org/10.1073/pnas.1407756111}.
Kvantově pravděpodobnostní modelování pořadí otázek.

\bibitem{ibm}
IBM Quantum Learning.
\emph{Mathematical formulations of measurements}.
\href{https://quantum.cloud.ibm.com/learning/en/courses/general-formulation-of-quantum-information/general-measurements/formulations-of-measurements}{Online: IBM Quantum Learning, kapitola o měření}.
Standardní měřicí formalismus.

\bibitem{pollock}
F.~A. Pollock, C. Rodríguez-Rosario, T. Frauenheim, M. Paternostro a K. Modi.
Operational Markov Condition for Quantum Processes.
\emph{Physical Review Letters} 120, 040405 (2018).
\url{https://doi.org/10.1103/PhysRevLett.120.040405}.
Operační popis časové paměti.

\bibitem{rac}
A. Ambainis, D. Leung, L. Mancinska a M. Ozols.
\emph{Quantum Random Access Codes with Shared Randomness}.
arXiv:0810.2937 (2008; revize 2009).
\url{https://arxiv.org/abs/0810.2937}.
Známý komunikační benchmark.

\bibitem{orch}
S. Hameroff a R. Penrose.
Consciousness in the universe: A review of the Orch OR theory.
\emph{Physics of Life Reviews} 11(1), 39--78 (2014).
\url{https://doi.org/10.1016/j.plrev.2013.08.002}.
Zdroj konkrétní biologicko-fyzikální hypotézy, nikoli její potvrzení.

\bibitem{repo}
J. Mikulík.
\emph{finite-Shannon-theorem}.
\url{https://github.com/jendamikulik/finite-Shannon-theorem}.
Navazující matematický projekt. Zde netvrdíme audit celého repozitáře ani jeho klinickou validaci.
\end{thebibliography}
\end{document}

